\HLASourcePage{18}{17}
\thispagestyle{hla-body}
The formal definition of the M\&S HLA comprises three main components: the HLA framework and rules
specification (found in IEEE Std 1516\texttrademark-2025), the HLA federate interface specification (found in this
document), and the HLA object model template (OMT) (found in IEEE Std 1516.2\texttrademark-2025). This standard,
IEEE Std 1516.1-2025, provides a complete description of the essential elements of the second component
of the HLA, the interface specification.\textsuperscript{\fontsize{7.98pt}{8pt}\selectfont 8}

The HLA is an integrated architecture that has been developed to provide a common architecture for M\&S.
The HLA requires that interfederate communications use a standard application programmer's interface
(API). This specification defines the standard services and interfaces to be used by the federates in order to
support efficient information exchange when participating in a distributed federation execution.
Additionally, the capability for reuse of individual federates that adhere to these standard services and
interfaces is increased.

This document provides a specification for the HLA functional interfaces between federates and the runtime
infrastructure (RTI). The RTI provides services to federates in a way that is analogous to how a distributed
operating system provides services to applications. These interfaces are arranged into seven basic service
groups, and an eighth group of support services, as follows:

\begin{hlaalphlist}
  \item Federation management
  \item Declaration management (DM)
  \item Object management
  \item Ownership management
  \item Time management
  \item Data distribution management (DMM)
  \item Support services
\end{hlaalphlist}

The eight service groups describe the interface between the federates and the RTI and between the federates
and the software services provided by the RTI for use by HLA federates. The RTI requires a set of services
from the federate that are referred to as ``RTI-initiated'' and are denoted with a \textdagger{} (printer's dagger) after the
service name.

\begingroup\setlength{\parskip}{0pt}
\vspace{0.2pt}
\HLAAnchor{clause:1.4.1}\pdfbookmark[2]{1.4.1 HLA federation object model (FOM) framework}{bookmark:clause:1.4.1}\HLASubsubclause{1.4.1 HLA federation object model (FOM) framework}
\vspace{9.4pt}
\endgroup

A concise and rigorous description of the object model framework is essential to the specification of the
interface between federates and the RTI and of the HLA services. The rules and terminology used to
describe a FOM are described in IEEE Std 1516.2-2025. A simulation object model (SOM) describes salient
characteristics of a federate to aid in its reuse and other activities focused on the details of its internal
operation. As such, a SOM is not the concern of the RTI and its services. A FOM, on the other hand, deals
with interfederate issues and is relevant to the use of the RTI. FOMs describe the set of object classes and
interaction classes chosen for a planned federation along with the attributes of the object classes and
parameters of the interaction classes.

Every HLA object instance is an instance of an object class found in the FOM. Object classes are chosen by
the object model designer to facilitate an appropriate organizational scheme. Each object class has a set of
attributes associated with it. An attribute is a distinct, identifiable portion of the object's state. In this
discussion, ``attribute designator'' refers to the attribute, and ``attribute value'' refers to its contents. From the
federation perspective, the set of all attribute values for an object instance shall completely define the state
of that instance. Federates may associate additional state information with an object instance that is not
communicated between federates, but this action is outside the purview of the HLA FOM.

\vspace{0.4pt}
\HLACompactFootnote{8}{Information on references can be found in Clause 2.}
\clearpage
